\chapter*{镜像对称概览}
\addcontentsline{toc}{chapter}{镜像对称概览}
\setcounter{chapter}{0}\setcounter{section}{0}\setcounter{equation}{0}
\SourcePage{7}
\Lecture{1}{2019年8月29日}
本课程属于代数几何专题课。我们将讨论某种形式的\Term{jingxiangduicheng}{镜像对称}，先从它的历史概览说起。
\section{计数镜像对称}
设 $X$ 是一个 \Term{calabiyau}{Calabi–Yau 流形}。我们将重点讨论 Calabi–Yau 三维流形。这意味着
$K_X=\det T_X^*\simeq\OO_X$，即典范丛作为全纯线丛是平凡的。通常，我们还要求 $b_1=0$，并且流形不可约。
\examnum{1}\begin{example}[（$\PP^4$ 中的五次超曲面）]\label{ex:0.1}
取齐次五次多项式 $f\in\CC[x_0,\ldots,x_4]$。如果 $f$ 足够一般，它在 $\PP^4$ 中的零点集就是光滑的，从而给出一个 Calabi–Yau 三维流形的例子。我们有
\Eq{0.1}{0\longrightarrow\II/\II^2\longrightarrow\Omega_{\PP^4}|_X\longrightarrow\Omega_X\longrightarrow0,}
其中 $\II=(f)\subset\OO_{\PP^4}$。此外，$\II/\II^2=\II\otimes_{\OO_{\PP^4}}\OO_X$ 是可逆层；上述序列的第一个映射把 $f$ 映为 $df$。因此
\Eq{0.2}{K_{\PP^4}|_X=\det\Omega_{\PP^4}|_X=(\II/\II^2)\otimes K_X,}
并且
\Eq{0.3}{K_{\PP^4}=\OO_{\PP^4}(-5).}
现在 $\II\hookrightarrow\OO_{\PP^4}$；把 $f$ 除以 $x_0^5$，可将其写成次数为零的有理表达式，其极点阶数为 $5$。因此，作为抽象层，
\Eq{0.4}{\II\simeq\OO_{\PP^4}(-5).}
于是
\Eq{0.5}{\II/\II^2\simeq\OO_X(-5).}
由张量积关系以及
\Eq{0.6}{K_{\PP^4}|_X\simeq\OO_X(-5),}
可知 $K_X\simeq\OO_X$，所以 $K_X$ 必须是平凡的。
\end{example}
现在，我们希望由此构造一种弦理论。这是一个非常微妙的过程。存在称为 $\mathrm{IIA}(X)$ 和 $\mathrm{IIB}(X)$ 的理论，它们正是镜像对称中相关的理论。这些理论来自以某个十维空间\footnote{这是为了使理论没有反常。}$\RR^{1,3}\times X$ 为靶空间的超对称 $\sigma$ 模型，也就是所谓的\Term{chaogongxingchanglun}{超共形场论} $\mathrm{SCFT}_A(X)$ 与 $\mathrm{SCFT}_B(X)$。它们是不同的理论，能够给出可观测量；例如，可以从这些理论计算 $X$ 的 Hodge 数。特别地，我们能计算 $h^{1,1}(X)$ 与 $h^{2,1}(X)$，它们分别对应某些物理变量。在 B 侧做同样的计算，得到的却是 $h^{2,1}(X)$ 与 $h^{1,1}(X)$。于是，我们设想还存在另一个 $X'$，
\SourcePage{8}
使这两个数不再被互换。更具体地说，我们观察到：对于某些非常特殊的 $X$，可以找到一个 Calabi–Yau 流形 $Y$，满足
\Eq{0.7}{\begin{aligned}\mathrm{SCFT}_A(X)&=\mathrm{SCFT}_B(Y),\\\mathrm{SCFT}_B(X)&=\mathrm{SCFT}_A(Y).\end{aligned}}
由此产生的想法是：如果这确实是弦理论的一个模型，那么我们应当能够交换
\Eq{0.8}{\begin{aligned}\mathrm{IIA}(X)&\mathrel{\text{“} =\text{”}}\mathrm{IIB}(Y),\\\mathrm{IIB}(X)&\mathrel{\text{“} =\text{”}}\mathrm{IIA}(Y).\end{aligned}}
不过，这个要求或许有些过强。
\subsection{拓扑扭曲}
Witten 在 1988 年引入了所谓的\Term{tuopuniuqu}{拓扑扭曲} $\sigma$ 模型。它给出一种完全不同的理论。我们得到两个理论，分别称为 $A(X)$ 与 $B(X)$。
\begin{Warning}{0.1}结果表明，$A(X)$ 最终计算的是 $\mathrm{IIB}(X)$ 理论在某些极限下的量。\end{Warning}
\begin{Note}{0.1}先验来看，它们与 $\mathrm{SCFT}_A(X)$ 和 $\mathrm{SCFT}_B(X)$ 并没有关系。\end{Note}
事实上，如果 $\mathrm{SCFT}_A(X)=\mathrm{SCFT}_B(Y)$，那么
\Eq{0.9}{\begin{aligned}A(X)&=B(Y),\\B(X)&=A(Y).\end{aligned}}
$A(X)$ 和 $B(X)$ 计算 $\mathrm{SCFT}_B(Y)$ 与 $\mathrm{SCFT}_A(X)$ 的某些极限，称为 \Term{yukawa}{Yukawa 耦合}。
注意，经过这一扭曲过程，$A(X)$ 只把 $(X,\omega)$ 看作辛流形，其中 $\omega$ 是 Kähler 形式；而 $B(X)$ 仅依赖于复流形 $(X,I)$。
\subsection{有用的计算}
人们对镜像对称如此兴奋，是因为它能够帮助我们完成此前无法完成的计算。
文献 \cite{ref4} 计算了五次超曲面及其镜像的 Yukawa 耦合。具体来说，他们计算了镜像五次超曲面 $Y_t$ 的 $F_B$，并断言它实际上等于五次超曲面的 $F_A$。从几何上看，$F_A$ 与亏格为零的全纯曲线的计数有关，而 $F_B$ 与\Term{zhouqijifen}{周期积分}有关：
\Eq{0.10}{\int_\alpha\Omega_{Y_t}=F_B(t),}
其中 $\alpha\in H_3(Y_t)$。因此，他们预测了某些曲线计数；后来有人直接算出了这些数目，结果彼此一致。
\begin{Warning}{0.2}$\Omega_{Y_t}$ 只在相差一个倍数的意义下确定，因此实际上应写成
\Eq{0.11}{F_B(t)=\exp\left(\frac{\int_{\alpha_1}\Omega_{Y_t}}{\int_{\alpha_0}\Omega_{Y_t}}\right),}
其中 $\alpha_i\in H_3(Y_t)$。
\end{Warning}
随后，Morrison 与 Deligne 于 1992 年用 Hodge 理论以及 Calabi–Yau 模空间的更多参数描述了 $F_B(Y)$。1993 年，Gromov–Witten 理论开始进入这一研究领域，使 $F_A(X)$ 获得精确的定义。这样，我们至少有了一个数学上的陈述。
\SourcePage{9}
\subsection{同调镜像对称}
1994 年，Kontsevich 在国际数学家大会上作了那场著名的报告，\Term{tongdiaojingxiangduicheng}{同调镜像对称}由此蓬勃发展。他指出：作为数学家，我们其实并不真正了解超共形场论；但真正应当成立的是
\Eq{0.12}{D\operatorname{Fuk}(X)=D^b(\OO_Y).}
这是一种表述，而不是一种解释。
Siebert 教授希望让我们相信，存在一种构造镜像对的方法。
\subsection{证明数值镜像对称}
1996 年，Givental 证明：在超曲面的情形下，$F_A$ 确实等于其镜像的 $F_B$。这一证明在某种意义上是通过计算表明两侧确实相等。Siebert 教授对此并不十分满意。1997 年，Lian、Liu 与 Yau 在更一般的情形下证明了这一结论。
\subsection{证明同调镜像对称}
2003 年，Paul Seidel 证明了 $\PP^3$ 中四次超曲面的同调镜像对称。大致而言，他证明两侧都有足够的刚性，因此只需进行很少的计算。Siebert 教授对此也并不十分满意。随后，N. Sheridan 于 2011 年证明了所有 Calabi–Yau 超曲面的这一结论。
\subsection{当前的发展状况}
镜像对称还有许多其他表现形式。事实上，甚至几何 Langlands 理论也可以视为某种形式的镜像对称。
在同调镜像对称方面，一些从事辛几何研究的人，正在尝试对所谓具有 SYZ 纤维化的辛流形证明这一结论，其镜像是一个刚性解析空间。
此外，还有内蕴构造。这是 Siebert 教授与 Mark Gross 合作研究过的方向，并且具有许多应用。其思想是把镜像对称用作数学中的工具，而不仅仅将其视为物理学中的一种现象。关键在于必须找到产生镜像的方法。
目前为止，这整个故事都发生在亏格为零的情形，也就是物理学家所说的树级。还存在高亏格的情形。从表示论的角度看，它与量子群有某种关系，称为二次量子化的镜像对称。\footnote{这一术语来自量子场论。}与之相关的还有一个完整的研究领域，称为拓扑递归。
\subsection{课程安排}
\begin{enumerate}
\item COGP 计算的一部分，即周期的计算。
\item Gromov–Witten 理论、虚拟基本类与模叠。
\item 环面退化与镜像构造。\footnote{这将包含一些环面几何的介绍。}
\item 证明同调镜像对称的一种策略，是计算两侧的齐次坐标环。Polishchuk 已证明，这个环决定 $D^b(\OO_Y)$。如果能在辛几何一侧进行相应的计算，就很理想，因为这将给出一种非常巧妙的同调镜像对称证明。
\item 高亏格情形：Donaldson–Thomas 不变量在镜像对称中扮演着某种尚不清楚的角色，因为它们将与高亏格的故事有关。可以用“晶体熔化”进行这些计算，这属于某种统计力学。
\end{enumerate}
